% Fragment partage : inclus par broadcast.tex (dossier autonome)
% et par dossier-conception-alanya.tex (dossier client assemble).
% Ne pas y mettre de preambule ni de \part : c'est l'appelant qui les pose.

\chapter{Ce que l'on ajoute}

\section{Le besoin}

Alanya n'a aujourd'hui aucun moyen de s'adresser à ses utilisateurs. Annoncer une
maintenance, une nouvelle version ou prévenir un commerçant que sa vérification expire
suppose de passer par un canal extérieur --- e-mail, réseaux sociaux --- que la plupart des
utilisateurs ne consultent pas.

On ajoute donc un \textbf{compte officiel} capable d'envoyer deux choses~:

\begin{itemize}[leftmargin=1.4em]
  \item des \textbf{messages}, qui apparaissent comme une conversation ordinaire~;
  \item des \textbf{statuts}, qui apparaissent dans le fil des statuts.
\end{itemize}

Dans les deux cas, à \textbf{tous} les utilisateurs ou à une \textbf{partie} d'entre eux.

\section{Ce qui a déjà été décidé}

Le socle de compte (étape~1) a posé les règles qui encadrent ce compte~:

\begin{center}
\begin{tabularx}{\linewidth}{@{}p{4.4cm}X@{}}
\toprule
\thd{Création} & Par un administrateur, jamais par une inscription ordinaire. \\
\thd{Apparence} & Avatar = logo Alanya, badge officiel doré, distinct du badge de
vérification. \\
\thd{Réponse} & \textbf{Impossible}. Le compte diffuse, il ne converse pas. \\
\thd{Nombre} & Le modèle en supporte plusieurs (support, actualités\ldots) sans
modification. \\
\bottomrule
\end{tabularx}
\end{center}

\chapter{Cibler une diffusion}

\section{Un critère, pas un choix parmi trois}

La spécification d'origine prévoyait trois cibles~: tous, par pays, ou une liste de comptes
choisis. C'est insuffisant --- on veut aussi pouvoir s'adresser aux commerçants, aux comptes
inscrits depuis plus de trois mois, à ceux vus récemment, et à d'autres critères qu'on ne
connaît pas encore.

Le ciblage devient donc une \textbf{combinaison de conditions}, que l'administrateur
assemble. Ajouter un critère demain ne demandera aucune modification de la base.

\begin{center}
\begin{tabularx}{\linewidth}{@{}p{4.6cm}X@{}}
\toprule
\thd{Pays} & Un ou plusieurs pays. \\
\thd{Genre de compte} & Personnel, business, officiel. \\
\thd{État de vérification} & Vérifié, en attente, expiré\ldots{} \\
\thd{Ancienneté} & Inscrit avant ou après une date, ou depuis une durée. \\
\thd{Activité} & Vu dans les $N$ derniers jours. \\
\thd{Comptes choisis} & Une liste explicite, pour un message de service. \\
\bottomrule
\end{tabularx}
\end{center}

\section{Savoir qui recevra, avant d'envoyer}

À chaque modification du ciblage, l'administrateur voit l'estimation se recalculer. C'est le
garde-fou principal~: on ne diffuse pas à trois cent mille personnes sans l'avoir vu écrit.

\begin{center}
  \imageOuCadre{maquettes/diffusion-composeur.png}{\linewidth}%
    {Composeur de diffusion : contenu, constructeur de critères, estimation en direct\\
     \texttt{maquettes/diffusion-composeur.png}}
\end{center}

\begin{note}[Maquette interactive]
Les conditions s'activent et se désactivent, l'estimation et la requête produite suivent~:\\
\href{https://claude.ai/code/artifact/b609a75d-ea57-490a-8df4-2fe31d9df609}%
{\texttt{claude.ai/code/artifact/b609a75d}}\\
Source versionnée~: \code{docs/architecture/maquettes/diffusion-admin.html}
\end{note}

\section{L'historique}

Chaque diffusion conserve son critère, et non sa liste de destinataires. Six mois plus tard,
on peut relire à qui elle s'adressait --- ce qu'une liste d'identifiants ne dirait pas.

L'historique affiche aussi le taux d'ouverture, dérivé du nombre de destinataires ayant
réellement rouvert l'application depuis l'envoi.

\begin{note}[La relance J-30 passe par ce canal]
Les avertissements d'expiration de vérification conçus à l'étape~1 ne sont rien d'autre
qu'une diffusion dont le critère est « business, vérifié, échéance dans moins de 30~jours ».
Aucun mécanisme spécifique n'est construit pour eux.
\end{note}

\chapter{Ce que voit l'utilisateur}

\section{Un message diffusé}

Une conversation ordinaire avec le compte Alanya, à ceci près que la zone de saisie est
remplacée par un bandeau~: \guillemotleft~Ce compte diffuse des annonces. Vous ne pouvez pas
y répondre.~\guillemotright{} Les actions \guillemotleft~répondre~\guillemotright{} et
\guillemotleft~réagir~\guillemotright{} n'apparaissent pas sur ses messages.

Tout le reste fonctionne normalement~: notification, compteur de non-lus, conservation hors
ligne, recherche. C'est voulu --- voir chapitre~\ref{ch:livraison}.

\section{Un statut diffusé}

Il apparaît dans le fil des statuts, avec l'avatar Alanya et le badge officiel, exactement
comme le statut d'un contact. Il expire au bout de 24~heures comme les autres.

\begin{alerte}[Un obstacle que la spécification d'origine n'avait pas vu]
Aujourd'hui, un statut n'est visible que si l'auteur et le lecteur sont \textbf{mutuellement}
en contacts préférés. En l'état, un statut officiel serait invisible pour absolument tout le
monde --- le compte Alanya ne peut pas ajouter trois cent mille personnes en favoris, ni
espérer qu'elles le fassent en retour. Le correctif tient en trois lignes et se trouve au
chapitre~\ref{ch:statuts-officiels}.
\end{alerte}
